% ===========================================================================
% Introduction and Related Work
% Paper: Topology-Aware Crack and Spalling Segmentation in Concrete Dams
%        via Dynamic Snake Convolution and Skeleton Recall Loss
% ===========================================================================

\section{Introduction}
\label{sec:intro}

Concrete dams are safety-critical infrastructure whose structural integrity must be continuously monitored throughout multi-decade service lives.
Surface defects---predominantly cracks and spalling---serve as early indicators of internal deterioration, yet their manual inspection remains labour-intensive, subjective, and logistically constrained by the scale and accessibility of dam faces~\cite{chen2024damsurvey,fang2024dam}.
Automated semantic segmentation of these defects from high-resolution imagery therefore promises substantial gains in inspection efficiency and repeatability.

Recent deep-learning methods have advanced pixel-level defect segmentation considerably.
Encoder--decoder architectures such as U-Net~\cite{ronneberger2015unet} and hierarchical transformers like SegFormer~\cite{xie2021segformer} achieve competitive Intersection-over-Union (IoU) on benchmark crack datasets~\cite{zou2019deepcrack,liu2019deepcrack}.
However, IoU and related region-overlap metrics are insensitive to topological errors: a prediction that fragments a continuous crack into disconnected segments may still score favourably on IoU while misrepresenting the structural extent of the defect~\cite{pantoja2024cracktopology}.
For structural safety assessment, the \emph{connectivity} of a predicted crack is as important as its pixel coverage---a fragmented prediction can lead an engineer to underestimate propagation length and thus misjudge residual capacity.

Topology-preserving losses such as clDice~\cite{shit2021cldice} and persistent-homology penalties~\cite{hu2019topology,clough2020topological} partially address this gap, yet they either incur prohibitive GPU memory for high-resolution dam images or operate on binary masks without multi-class support.
The recently proposed Skeleton Recall Loss (SRL)~\cite{shit2024skeleton} overcomes both limitations by computing recall on CPU-extracted ground-truth skeletons, achieving strong connectivity preservation with over 90\% reduction in computational overhead relative to clDice.
Concurrently, Dynamic Snake Convolution (DSConv)~\cite{qi2023dynamic} reshapes kernel sampling grids to follow elongated, tortuous paths---a morphological prior that aligns naturally with the geometry of cracks.

In this paper we propose \textbf{DSCFormer}, a dual-branch architecture that augments a SegFormer-B2 backbone with a parallel DSConv feature stream and trains with a composite loss incorporating SRL.
We evaluate on a new \textbf{DamSegment} benchmark of 1\,500 annotated images spanning cracks and spalling on concrete dam surfaces.
Our contributions are:
\begin{itemize}
    \item A topology-aware segmentation framework combining dynamic snake convolution with skeleton-based supervision for fine-grained dam defect segmentation, achieving 72.3\% mIoU$_{fg}$ with test-time augmentation.
    \item Integration of Skeleton Recall Loss into a multi-class setting, demonstrating that explicit skeleton supervision yields +4.8 clDice and +5.3 Boundary F1 (BF1) over the SegFormer-B2 baseline without increasing inference cost.
    \item A systematic comparison with foundation model adaptation approaches (SAM-LoRA, DINOv2-LoRA), revealing that task-specific architectural design outperforms parameter-efficient fine-tuning of large vision models on small-scale industrial datasets.
    \item The DamSegment dataset with pixel-level crack and spalling annotations, providing the first public benchmark specifically targeting concrete dam surface defects.
\end{itemize}

\section{Related Work}
\label{sec:related}

\subsection{Concrete Defect Segmentation}
\label{sec:related:defect}

Early CNN-based crack detectors such as CrackNet treated defect identification as patch classification~\cite{zou2019deepcrack}.
Fully convolutional designs subsequently enabled dense pixel-wise prediction: DeepCrack~\cite{liu2019deepcrack} introduced deeply-supervised side outputs fused across encoder--decoder stages, while Zou~\emph{et~al.}~\cite{zou2019deepcrack} exploited hierarchical SegNet features with multi-scale fusion.
These methods established the encoder--decoder paradigm that remains dominant in structural defect segmentation.

More recently, transformer-based architectures have demonstrated improved global context modelling.
CrackFormer~\cite{liu2021crackformer} embeds self-attention within a SegNet-like structure to capture long-range crack continuity, and SegFormer~\cite{xie2021segformer} provides a lightweight hierarchical transformer encoder paired with an all-MLP decoder that achieves strong accuracy--efficiency trade-offs.
Universal architectures such as Mask2Former~\cite{cheng2022mask2former} and hierarchical vision transformers like Swin Transformer~\cite{liu2021swin} further push segmentation accuracy on general benchmarks but have seen limited adoption for infrastructure defects due to annotation scarcity and domain shift.

Applied to dam inspection specifically, Arafin~\emph{et~al.}~\cite{arafin2024concrete} performed multi-class concrete defect segmentation using DeepLabV3+, while Fang and Liu~\cite{fang2024dam} proposed a lightweight MobileNetV2-based architecture for microscopic dam crack quantification.
Chen~\emph{et~al.}~\cite{chen2024damsurvey} surveyed visual perception methods for intelligent dam hub inspection, identifying limited dataset scale and poor topology preservation as key open challenges.

Despite progress, these methods optimise pixel-level overlap metrics (IoU, Dice) and do not explicitly penalise topological disconnections---a critical shortcoming for thin, elongated cracks whose connectivity carries diagnostic significance.

\subsection{Topology-Aware Segmentation}
\label{sec:related:topology}

Topology-aware approaches seek to enforce structural constraints beyond per-pixel accuracy.
Hu~\emph{et~al.}~\cite{hu2019topology} first introduced a differentiable topological loss based on persistent homology, penalising Betti-number discrepancies between predicted and ground-truth segmentations.
Clough~\emph{et~al.}~\cite{clough2020topological} extended persistent-homology supervision to cardiac imaging, demonstrating that topological correctness and pixel accuracy are complementary objectives.
However, persistent-homology computation scales poorly to high-resolution images common in infrastructure inspection.

For tubular structures, Shit~\emph{et~al.}~\cite{shit2021cldice} proposed clDice, a differentiable soft-skeleton metric that guarantees topology preservation up to homotopy equivalence for binary masks.
While elegant, clDice requires GPU-based iterative morphological skeletonisation during training, incurring substantial memory overhead at full resolution and lacking native multi-class support.
The Skeleton Recall Loss~\cite{shit2024skeleton} addresses both limitations: it pre-computes ground-truth skeletons on CPU and formulates a recall objective over skeleton voxels, achieving superior connectivity metrics on five tubular-structure benchmarks with over 90\% reduction in GPU memory relative to clDice.

Recent work on crack-specific topology includes the Crack Topology Score (CTS)~\cite{pantoja2024cracktopology}, a skeleton-matching evaluation protocol weighted by crack length, and boundary-guidance methods~\cite{wang2024dsctopology} that incorporate edge priors into encoder--decoder training.
These advances highlight a growing consensus that pixel-level metrics alone are insufficient for safety-critical crack assessment.

Our work builds upon SRL~\cite{shit2024skeleton} by extending it to the multi-class dam defect setting and pairing it with a geometry-aware feature extractor, yielding complementary gains in both connectivity and boundary accuracy.

\subsection{Deformable and Dynamic Convolutions for Elongated Structures}
\label{sec:related:deformable}

Standard convolutions sample on fixed rectangular grids, an inductive bias poorly matched to thin, curvilinear structures.
Deformable Convolutional Networks (DCN)~\cite{dai2017deformable} introduced learnable offsets that adaptively shift sampling locations, enabling feature extraction along object boundaries.
DCNv2~\cite{zhu2019deformable} augmented offsets with per-sample modulation amplitudes, improving the network's ability to suppress irrelevant context in cluttered scenes.

While DCN offsets are unconstrained and learned purely from data, the Dynamic Snake Convolution (DSConv)~\cite{qi2023dynamic} imposes a \emph{topological geometric constraint}: offsets are parameterised along a one-dimensional path, forcing the kernel to deform into a snake-like shape that naturally traces elongated, tortuous structures such as vessels and cracks.
DSConv demonstrated state-of-the-art vessel segmentation by combining this morphological prior with multi-view feature fusion.

Pyramid pooling~\cite{zhao2017pspnet} and multi-scale attention mechanisms~\cite{liu2021swin} complement deformable operators by providing broader context, yet they do not explicitly encode the elongated geometry of thin defects.
The concurrent DSCformer architecture applies enhanced DSConv to pavement crack segmentation, confirming the effectiveness of snake-shaped receptive fields for this domain.

Our DSCFormer differs from prior DSConv applications in two respects: (i)~we integrate the DSConv branch as a \emph{parallel} stream alongside a SegFormer-B2 encoder rather than replacing standard convolutions wholesale, preserving the transformer's global context; and (ii)~we couple the geometric inductive bias of DSConv with the topological supervision of Skeleton Recall Loss, yielding synergistic improvements in both boundary fidelity and connectivity.
